% !PRINTED 54
\weberpage{54}
%<<CONT

stellungsweise doch nicht ganz gleichgültig ist, die erste oft auch den Vorzug verdient. Besonders in den Fällen, wo die Coëfficienten Zahlen sind und es auf das zahlentheoretische Verhalten dieser Coëfficienten ankommt, darf man nicht ausser Acht lassen, dass durch die Binomialcoefficienten ein der Sache fremdes numerisches Element eingeführt wird. Dass Gauss in der Theorie der quadratischen Formen (in den Disq. ar.) die Schreibweise mit den Binomialcoefficienten anwendet, wenn er die quadratischen Formen durch $ax^2+2bxy+cy^2$ darstellt, und dass diese Bezeichnung allgemein Eingang gefunden hat, hat in der Zahlentheorie zu einer unnöthigen und sehr bedauerlichen Complication geführt.

\begin{center}
{\large \S.\ 14.}

\vspace{0.5em}
{\large\gesperrt{Derivirte eines Productes.}}
\end{center}

Die derivirten Functionen, die wir hier betrachtet haben, sind keine anderen als die aus der Differentialrechnung bekannten Differentialquotienten; wir haben den Begriff aber hier, wo es sich um ganze rationale Functionen handelt, ohne Anwendung der Infinitesimalrechnung gewonnen aus den Entwickelungscoefficienten der Potenzen von $y$ in der Function $f(x+y)$. Bezeichnen wir die $\nu$te Ableitung von $f(x)$ mit $D_\nu f$, so ist nach (2), \S~13:
\begin{equation}
 f(x+y)=f(x)+yD_1f+{y^2\over1\cdot2}D_2f+{y^3\over1\cdot2\cdot3}D_3f+\cdots,
\tag{1}
\end{equation}
und daraus ergeben sich sofort die beiden Grundsätze, die sich in den Formeln
\begin{equation}
 D_\nu(Cf)=C D_\nu f;
\tag{2}
\end{equation}
\begin{equation}
 D_\nu(f+\varphi)=D_\nu f+D_\nu\varphi
\tag{3}
\end{equation}
ausdrücken, worin $C$ eine Constante, $\varphi$ eine zweite ganze rationale Function von $x$ ist.

Eine Verallgemeinerung der Formel (2) giebt die Darstellung der Derivirten des Productes $f\varphi$. Setzt man nämlich nach (1) abkürzend
\begin{align}
 f(x+y)&=u_0+yu_1+y^2u_2+\cdots+y^n u_n,\notag\\
 \varphi(x+y)&=v_0+yv_1+y^2v_2+\cdots+y^m v_m,
\tag{4}
\end{align}
also
\begin{equation}
 u_\nu={D_\nu f\over\Pi(\nu)},\qquad
 v_\nu={D_\nu\varphi\over\Pi(\nu)},
\tag{5}
\end{equation}
%CONT>>
